% !TEX root = ../main.tex

\section{Introduction}
\label{ch:introduction}

What the paper is about. Motivation for the problem at hand. Always in your own words, your understanding.

Recently, Transformers have revolutionized the machine learning field across various domains, including natural language processing (NLP) and computer vision. Prominently, large Transformer-based language models, or LLMs, like GPT-3 by Brown et al. in 2020, have shown remarkable general-purpose capabilities.These models are now widely integrated into user-facing applications, including chatbots, search engines, and code assistants. Large language models (LLMs) such as GPT4 (OpenAI, 2023), LLaMA-2 (Touvron et al.,2023), Claude-3 (AnthropicAI, 2023), and Gemini (Team et al., 2023) have led to tremendous advancements in language understanding and generation, achieving state-of-the-art performance in a wide array of real-world tasks.  In computer vision, Vision Transformers (ViTs) have demonstrated superior performance in image classification and object detection (Dosovitskiy et al., 2020). Additionally, Transformers are being employed in scientific domains for tasks such as protein folding prediction with models like AlphaFold  More recently, Transformer models have been applied to a variety of scientific tasks ranging from foundation models in fluid dynamics to astrophysics and climate science (Nguyen et al., 2023; McCabe et al., 2023; Lanusseet al., 2023). However, despite their impressive capabilities, all these models are like a black box which takes some input and prodice some output. We do not know how the model process each part of its input and how it arrives to its decision. Understanding these models is difficult because these models consist of densely connected layers of non linear complex functions.

To understand its working, many researchers try to analyse weights of attention weights and hidden representations. In What does BERT look at? An Analysis of BERT’s Attention, the author try to understand what is the function of each head. Another line of work investigates the internal vector representations of the model (Adi et al.,2017; Giulianelli et al., 2018; Zhang and Bowman, 2018,BERT Rediscovers the Classical NLP Pipeline), often using probing classifiers. Probing classifiers are simple neural networks that take the internal vector representations of a pre-trained model as input. They are trained to do a supervised task (e.g., part-of-speech tagging). If a probing classifier achieves high accuracy, it suggests that the vector representations reflect the corresponding aspect of language(e.g., low-level syntax). These methods provide partial insight intomodel behavior but have also been shown to be misleading [Jain and Wallace, 2019, Bolukbasi et al., 2021]; they do not tell algorithically what the model is doing.

